\section{Introduction}
\label{sec:intro}

Retrieval-augmented generation answers a query by re-reading raw source chunks at serving time. This is simple and, because it never modifies what it read, hard to corrupt: every answer is grounded in whatever the retriever pulled back at that moment. It is also wasteful and structurally blind to change --- if a fact has been superseded, RAG has no mechanism to prefer the new version over the old one unless the retriever happens to rank recency, and most do not.

An alternative pattern, popularized by Karpathy as an ``LLM-Wiki'' \citep{karpathy2026wiki}, compiles a corpus once (and re-compiles or incrementally updates it as new documents arrive) into a curated, LLM-written store: a wiki that an LLM reads and rewrites, rather than a chunk index an LLM merely searches. Compiling gives a system the chance to resolve contradictions once, at write time, instead of re-litigating them on every query. That chance is also the risk: a compiled store that gets this wrong does not just retrieve a stale chunk, it \emph{writes a stale answer into the store itself}, and every subsequent query inherits the error.

This paper asks what a compiled store needs in order for that promise to be trustworthy rather than merely convenient. We organize the answer as a four-stage lifecycle, each stage validated by an independent line of experiments:

\begin{itemize}[leftmargin=1.4em,itemsep=0.2em]
\item \textbf{Compile} (\S\ref{sec:compile}): does compiling actually fix the failure mode RAG has --- answering with a superseded fact --- and is the fix coming from having a wiki (structure, memory) or from something narrower?
\item \textbf{Certify} (\S\ref{sec:certify}): once compiled, how does a store's owner or consumer know, with a stated confidence level rather than a point estimate, how much of it is still right?
\item \textbf{Maintain} (\S\ref{sec:maintain}): as new documents force the store to be rewritten, how much of what it already knew survives the rewrite, and can maintenance be made both cheap and hard to game?
\item \textbf{Retract} (\S\ref{sec:retract}): when a source is withdrawn or a fact must be removed for a real reason (legal, safety, correction), does removing it from the store actually make it stay gone?
\end{itemize}

Each stage is validated with its own benchmark and real systems, and in one case cross-checked against an independent third-party tool for external validity. This paper has not been externally peer-reviewed; every result was subjected to internal scrutiny and is reported only to the extent it survives that scrutiny. Rather than present the strongest-sounding number for each stage, we state exactly what statistical support each claim has, and are explicit in \S\ref{sec:limitations} about what is asserted but not yet run. We believe this is the more useful contribution: a lifecycle view of what actually needs to be true for a compiled knowledge store to be safe to depend on, grounded in real numbers with their real caveats attached, rather than a single headline result.

\textbf{Contributions.} (1) A controlled dissociation showing that recency-aware \emph{resolution}, not graph structure or persistent memory, is what eliminates staleness in compiled stores, validated against two real structured RAG systems (\S\ref{sec:compile}). (2) A finite-sample, noisy-judge-corrected statistical contract for compiled-store fidelity and currency, maintainable incrementally at a fraction of full-audit cost, which surfaces and quantifies two failure modes --- build-to-build disagreement and self-audit overclaiming --- that a point-estimate evaluation would miss entirely (\S\ref{sec:certify}). (3) Evidence that faithful maintenance under repeated rewriting is partially trainable at equal deploy cost, wrapped in a per-batch anytime-valid retention bound that catches and rolls back harmful updates rather than only measuring damage after the fact (\S\ref{sec:maintain}). (4) A demonstration that retraction from an agent's compiled memory is not durable under the agent's own consolidation loop unless the deletion gate is keyed on membership rather than correctness, with a fix that reduces recoverability from statistically-indistinguishable-from-no-deletion to matching a never-ingested oracle (\S\ref{sec:retract}).
